\chapter{Margin}\label{ch:m1:margin}

A fund goes home on Friday long four hundred index futures, fully margined.
On Monday it holds the same four hundred contracts and owes the clearing
house cash twice over: once for what the contracts lost, and again because
the clearing house has decided, looking at the same prices, that each
contract is now more dangerous to hold. Over March 2020 the \omterm{def:m1:clearing-and-settlement:margin}{initial margin}
required by central counterparties worldwide rose by roughly \$300 billion,
and on a single day, 9 March, they called \$140 billion of \omterm{def:m1:clearing-and-settlement:margin}{variation margin}.
None of that money was a loss for the system as a whole; all of it had to be
found in cash within hours by someone. \Cref{ch:m1:clearing-and-settlement}
said what margin is for. This chapter computes it, the way the clearing
house does, and looks at how the computation behaves on the day it matters.

\section{Variation margin}

\omterm{def:m1:clearing-and-settlement:margin}{Variation margin} (\cref{def:m1:clearing-and-settlement:margin}) is
arithmetic: position times the change in the settlement price times the
multiplier, paid in cash, every day, in the contract's currency. Three
properties make it harder than it looks. It is \emph{asymmetric in
liquidity}: a hedged portfolio whose other leg is not at the same clearing
house pays cash on the losing leg and receives none on the winning one. It is
\emph{fast}: a call made in the morning is due the same morning, and in
stressed markets clearing houses call again during the day. And it is
\emph{unforgiving}: a member that does not pay is in default.

\begin{definition}[Performance bond and maintenance margin]\label{def:m1:margin:bond}
A \emph{performance bond}\index{performance bond} is the futures industry's
name for \omterm{def:m1:clearing-and-settlement:margin}{initial margin}: collateral deposited to guarantee performance, not a
down payment. For a client account the broker sets an \emph{initial} level,
required to open a position, and a lower \emph{maintenance
margin}\index{maintenance margin}: when the account's equity falls below the
maintenance requirement the client is called to restore it to the initial
level. For clearing members the two coincide.
\end{definition}

\begin{definition}[Intraday margin call]\label{def:m1:margin:intraday}
An \emph{intraday margin call}\index{intraday margin call} is a call for
variation or \omterm{def:m1:clearing-and-settlement:margin}{initial margin} made by a clearing house between its scheduled
daily cycles, on the basis of current prices and positions, payable within
the hour or so that its rules specify.
\end{definition}

\section{Initial margin by scenarios}

\begin{definition}[Scanning range and risk array]\label{def:m1:margin:scan}
In a scenario-based margin system the clearing house sets, for each product,
a \emph{price scan range} and a \emph{volatility scan range}: the largest
moves in price and in implied volatility it wishes to cover over its
horizon, the \emph{scanning range}\index{scanning range}. The \emph{risk
array} of a contract is its gain or loss under each of a fixed list of
scenarios built from those ranges; the \emph{scan risk} of a portfolio is its
largest total loss over the list.
\end{definition}

The system that carried futures margining for three decades, SPAN, uses
sixteen scenarios. Two move volatility alone, up and down. Twelve combine a
price move of one third, two thirds or the whole of the price scan range, up
or down, with volatility up or down. The last two are \emph{extreme} moves, a
multiple of the scan range, of which only a fraction of the loss is counted:
they exist for short options far out of the money, which lose nothing in the
first fourteen. For a future, the price scan range is the \omterm{def:m1:margin:bond}{maintenance margin}
itself.

\begin{example}[The risk array of one future]\label{ex:m1:margin:array}
A regulator's review of the system prints the array of the large S\&P 500
future on 12 April 2001, when its \omterm{def:m1:margin:bond}{maintenance margin} was \$17\,250: zero in
the two volatility scenarios; losses of $\mp\$5\,750$, $\mp\$11\,500$ and
$\mp\$17\,250$ at one, two and three thirds; and in the extreme scenarios, a
move of three times the range with 30\% of the loss counted, $\pm\$15\,525$. A
long future's \omterm{def:m1:margin:scan}{scan risk} is the full move down, \$17\,250; the extreme
scenario, at \$15\,525, does not bind. It binds for the option seller.
\end{example}

\begin{method}[A scenario margin]\label{met:m1:margin:span}
For each product:
\begin{enumerate}
\item \textbf{\omterm{def:m1:margin:scan}{Scan risk}}: sum position times risk array over all contracts,
scenario by scenario; take the largest loss.
\item \textbf{Inter-month charge}: futures of different expiries offset
fully in the scan; add a charge per spread for the risk that they do not
move together.
\item \textbf{Short option minimum}: a floor per short option, compared with
the sum of the first two.
\end{enumerate}
Across products: subtract an \emph{inter-commodity spread
credit}\index{inter-commodity spread credit} for pairs whose prices are
correlated and whose positions offset. Add the results.
\end{method}

\begin{omfigure}
\begin{tikzpicture}
\begin{axis}[width=11.5cm,height=5cm,ybar,bar width=8pt,
  xlabel={scenario},ylabel={portfolio loss (\$ thousand)},
  tick label style={font=\small},label style={font=\small},
  xmin=0.3,xmax=16.7,xtick={1,2,3,4,5,6,7,8,9,10,11,12,13,14,15,16},ymin=-50,ymax=140,ymajorgrids,grid style={gray!25},table/col sep=comma]
\addplot[fill=omDef!60,draw=omDef] table[x=scenario,y=loss_k]{figdata/markets-1/20-margin/scenarios.csv};
\node[font=\footnotesize,anchor=south east] at (axis cs:16.6,112) {scan risk: 107.8};
\end{axis}
\end{tikzpicture}

\omcaption{Sixteen scenarios for a portfolio long 2 futures at 6\,000, short
10 puts struck at 5\,600 and short 10 calls at 6\,400 (three months, 18\%
volatility, price scan 345 points, volatility scan 4 points, extreme moves
of three ranges counted at 30\%). Odd scenarios raise volatility, even ones
lower it. The margin is set by scenario 16, the extreme fall. Data: the
chapter's build.}\label{fig:m1:margin:scenarios}
\end{omfigure}

\begin{omfigure}
\begin{tikzpicture}
\begin{axis}[width=11cm,height=5.2cm,
  xlabel={move in the future (points)},ylabel={P\&L (\$ thousand)},
  tick label style={font=\small},label style={font=\small},
  xmin=-1100,xmax=1100,xtick={-1000,-500,0,500,1000},ymin=-420,ymax=80,grid=major,grid style={gray!25},
  legend columns=3,legend style={font=\footnotesize,at={(0.5,-0.3)},anchor=north,draw=none,fill=none,/tikz/every even column/.append style={column sep=8pt}},
  table/col sep=comma]
\addplot[omThm,very thick] table[x=move,y=vol_up_k]{figdata/markets-1/20-margin/profile.csv};
\addplot[omDef,very thick,dashed] table[x=move,y=vol_down_k]{figdata/markets-1/20-margin/profile.csv};
\addplot[only marks,mark=*,mark size=1.8pt,black] table[x=move,y=pnl_k]{figdata/markets-1/20-margin/points.csv};
\draw[thin] (axis cs:-345,-420) -- (axis cs:-345,80);
\draw[thin] (axis cs:345,-420) -- (axis cs:345,80);
\legend{volatility $+4$ points,volatility $-4$ points,scenarios 1 to 14}
\end{axis}
\end{tikzpicture}

\omcaption{What the scan sees. The same portfolio revalued over a wide range
of prices; the fourteen standard scenarios are the dots, all inside the price
scan range (thin lines). The losses that matter for this portfolio lie
outside it, which is what scenarios 15 and 16 are
for.}\label{fig:m1:margin:profile}
\end{omfigure}

\begin{definition}[Portfolio margining]\label{def:m1:margin:portfolio}
\emph{Portfolio margining}\index{portfolio margining} sets the requirement on
the net risk of all positions in an account, recognising offsets between
them, in contrast with a sum of requirements computed position by
position.
\end{definition}

\begin{omfigure}
\begin{tikzpicture}
\begin{axis}[width=9.5cm,height=4.4cm,ybar,bar width=16pt,
  ylabel={margin (\$ thousand)},
  symbolic x coords={future,short put,short call,portfolio},xtick=data,
  tick label style={font=\small},label style={font=\small},xticklabel style={font=\small\strut},
  ymin=0,ymax=135,ymajorgrids,grid style={gray!25},enlarge x limits=0.18,
  nodes near coords,nodes near coords style={font=\footnotesize,/pgf/number format/.cd,fixed,fixed zerofill,precision=1},
  table/col sep=comma]
\addplot[fill=omDef!60,draw=omDef] table[x=item,y=margin_k]{figdata/markets-1/20-margin/offsets.csv};
\end{axis}
\end{tikzpicture}

\omcaption{Offsets. Margined separately the three legs require \$213\,000;
together, \$108\,000, because the put and the call cannot both lose. Data:
the chapter's build.}\label{fig:m1:margin:offsets}
\end{omfigure}

\section{Initial margin by value-at-risk}

Sixteen scenarios per product and a table of credits between products are
transparent and replicable, which is why the method lasted. They are also
crude: the offsets are set by committee, and a portfolio of many products is
margined as a sum of pairs. The large clearing houses have moved, or are
moving, to portfolio value-at-risk.

\begin{definition}[Margin period of risk]\label{def:m1:margin:mpor}
The \emph{margin period of risk}\index{margin period of risk} is the time
assumed between a member's last payment of \omterm{def:m1:clearing-and-settlement:margin}{variation margin} and the moment
its portfolio has been hedged or liquidated by the clearing house: the
horizon of the initial-margin calculation. It is a day or two for liquid
listed futures and longer for cleared swaps and illiquid products.
\end{definition}

\begin{dated}{2026-09}{From scenarios to simulation}\label{dat:m1:margin:models}
The Chicago clearing house's successor to SPAN is built on historical
value-at-risk: thousands of scenarios from at least ten years of data, with
volatility and correlation scaling, and separate reporting of market,
liquidity and concentration components. The Frankfurt clearing house's
portfolio method uses filtered historical simulation together with a
stress-period value-at-risk, computed by groups of products that would be
liquidated together, each with its own \omterm{def:m1:margin:mpor}{margin period of risk}.
\end{dated}

\begin{method}[Filtered historical simulation]\label{met:m1:margin:fhs}
(i) Estimate each day's volatility $\sigma_t$ from past returns, for instance
by exponential weighting. (ii) Standardise the last $n$ returns:
$z_s = r_s/\sigma_s$. (iii) Rescale to today: $\tilde r_s = z_s\,\sigma_t$.
(iv) Revalue the portfolio under each $\tilde r_s$, take the loss quantile at
the chosen confidence, and scale to the \omterm{def:m1:margin:mpor}{margin period of risk}.
\end{method}

Filtering makes the margin react within days to a change of regime, where a
plain historical window reacts only as extreme days accumulate. It is the
right property for covering tomorrow's loss, and the wrong one for the
stability of the system.

\section{Procyclicality}

\begin{definition}[Procyclicality]\label{def:m1:margin:procyclical}
A margin system is \emph{procyclical}\index{procyclicality} to the extent
that its requirements rise in stressed markets and fall in calm ones,
adding to the demand for liquidity exactly when liquidity is scarce.
\end{definition}

\begin{omfigure}
\begin{tikzpicture}
\begin{axis}[width=11cm,height=5.2cm,
  xlabel={trading days from the start of the stress},ylabel={margin (\% of position)},
  tick label style={font=\small},label style={font=\small},
  xmin=-100,xmax=299,ymin=0,ymax=10.5,grid=major,grid style={gray!25},
  legend columns=3,legend style={font=\footnotesize,at={(0.5,-0.3)},anchor=north,draw=none,fill=none,/tikz/every even column/.append style={column sep=8pt}},
  table/col sep=comma]
\addplot[omExo,very thick] table[x=day,y=hs]{figdata/markets-1/20-margin/procyclical.csv};
\addplot[omThm,very thick] table[x=day,y=fhs]{figdata/markets-1/20-margin/procyclical.csv};
\addplot[omDef,very thick,dashed] table[x=day,y=floored]{figdata/markets-1/20-margin/procyclical.csv};
\legend{historical simulation,filtered,filtered with a 25\% stressed floor}
\end{axis}
\end{tikzpicture}

\omcaption{Margin on a fixed position while daily volatility jumps from
0.8\% to 3.5\% for 25 days and decays (99\%, two-day horizon, 250-day
window). Over the worst ten days the plain margin rises by 109\%, the
filtered margin by 203\%, the floored margin by 94\%. The floor is paid for
in calm times: 4.5\% of the position against 3.0\%. Data: the tutorial's
simulation.}\label{fig:m1:margin:procyclical}
\end{omfigure}

The joint review of March 2020 by the banking, market-infrastructure and
securities standard setters found the rise in cleared \omterm{def:m1:clearing-and-settlement:margin}{initial margin} large
and uneven across asset classes and clearing houses, and listed among its
areas for further work the transparency of cleared margin, the liquidity
preparedness of market participants, and the responsiveness of
initial-margin models to stress. For a trading firm the operational
conclusion does not wait for policy: the liquidity needed to hold a position
is its \omterm{def:m1:clearing-and-settlement:margin}{variation margin} in a bad week \emph{plus} the increase of its \omterm{def:m1:clearing-and-settlement:margin}{initial
margin} in that same week, and both must be modelled.

\section{Tutorial: a scenario margin for futures and options}\label{tut:m1:margin:span}

\begin{tutorial}
\textbf{Goal.} Build risk arrays for futures and options on futures, scan a
portfolio, add the charges, and compare a filtered with an unfiltered
value-at-risk margin through a volatility shock. \textbf{End state:} the
four data figures of this chapter.
\begin{tutsteps}
\item \textbf{Sixteen scenarios} from two ranges and two extreme-move
parameters.
\omcode{code/firm/margin/firm_margin.py}{53}{60}{The scenario list: price move, volatility move, fraction of the loss counted.}\label{lst:m1:margin:scenarios}
\item \textbf{A risk array} by full revaluation with the Black (1976)
formula for options on futures.
\omcode{code/firm/margin/firm_margin.py}{63}{74}{Revaluation and the risk array of one long unit.}\label{lst:m1:margin:array}
\item \textbf{Scan, charge, floor.}
\omcode{code/firm/margin/firm_margin.py}{86}{102}{The margin of one product's positions.}\label{lst:m1:margin:product}
\item \textbf{\omterm{def:m1:margin:procyclical}{Procyclicality}.} Rescale the window to today's volatility.
\omcode{code/markets-1/20-margin/python/margin_demo.py}{30}{33}{Filtered historical simulation.}\label{lst:m1:margin:fhs}
\end{tutsteps}
\textbf{What to change next.} Replace the two extreme scenarios by a full
grid out to five ranges and see which portfolios were under-margined. Then
give the filtered method a floor on $\sigma_t$ instead of a stressed add-on
and compare the two tools at equal \omterm{def:m1:pnl-and-positions:realised}{average cost}.
\end{tutorial}

\section{Build: the margin engine}\label{bld:m1:margin:engine}

\begin{build}
\bfield{Purpose} The miniature firm's clearing component
(\cref{ch:m1:clearing-and-settlement}) needs a number every night and on
demand: the \omterm{def:m1:clearing-and-settlement:margin}{initial margin} of each account, recomputable by the member.
\bfield{Interface} \texttt{Params(price\_scan, vol\_scan, extreme\_mult,
extreme\_cover, intermonth\_charge, short\_option\_min, multiplier)};
\texttt{Position(root, expiry, qty, strike, right)}; \texttt{Market};
\texttt{scenarios(params)}; \texttt{risk\_array(position, market, params)};
\texttt{product\_margin(positions, market, params)} returning \omterm{def:m1:margin:scan}{scan risk}, the
binding scenario, the inter-month charge, the short option minimum and the
total; \texttt{portfolio\_margin(by\_root, market, params, credits)}.
\bfield{Rules} Contract data from the master of
\cref{bld:m1:futures-contracts-and-exchanges:master}. Full revaluation, no
delta approximation. Futures of different expiries net in the scan and pay
the inter-month charge per spread. The inter-commodity credit is given only
when the two products' binding scenarios are price moves in opposite
directions.
\bfield{Acceptance tests} \texttt{code/firm/margin/tests/}: the future's
array of \cref{ex:m1:margin:array}, to the dollar; a \omterm{def:m1:matching-algorithms-and-implied-spreads:calendar}{calendar spread}; a short
far out-of-the-money put caught by the extreme scenario, and one caught by
the minimum; a covered option; credits granted and refused.
\bfield{Stretch} A second engine, portfolio value-at-risk by filtered
historical simulation, behind the same interface, and a report of the
difference between the two on every account.
\end{build}

\begin{omsources}
\item US Commodity Futures Trading Commission, \emph{Review of Standard
Portfolio Analysis of Risk (``SPAN'') Margin System}, April 2001.
\item Basel Committee on Banking Supervision, Committee on Payments and
Market Infrastructures and IOSCO, \emph{Review of margining practices},
final report and press release, 29 September 2022.
\item B.~H.~Cohen and K.~Tracol, ``Market turbulence and soaring margins:
lessons from two recent episodes'', \emph{BIS Quarterly Review}, March 2023.
\item CME Group, \emph{SPAN 2 methodology}; Eurex Clearing, \emph{Eurex
Clearing Prisma}.
\item F.~Black, ``The pricing of commodity contracts'', \emph{Journal of
Financial Economics} 3 (1976).
\end{omsources}

\section{Exercises}

\begin{exercise}[$\star$]\label{exo:m1:margin:1}
An account is long 25 E-minis. The settlement price goes from 6\,012.50 to
5\,968.25. Give the \omterm{def:m1:clearing-and-settlement:margin}{variation margin}, with its direction.
\end{exercise}

\begin{exercise}[$\star$]\label{exo:m1:margin:2}
A broker requires \$23\,100 of initial and \$21\,000 of \omterm{def:m1:margin:bond}{maintenance margin} per
contract. A client holds 5 contracts with \$120\,000 of equity and loses
\$3\,200 per contract. Is there a call, and for how much?
\end{exercise}

\begin{exercise}[$\star$]\label{exo:m1:margin:3}
Write the sixteen-entry risk array of one long future with a multiplier of
\$50, a price scan range of 345 points, and extreme moves of three ranges
counted at 30\%. Which entry is its margin?
\end{exercise}

\begin{exercise}[$\star\star$]\label{exo:m1:margin:4}
In \cref{fig:m1:margin:scenarios}, why do the odd scenarios lose more than
their even neighbours? Why is scenario 16 worse than scenario 15 although the
strangle is symmetric? What would make scenario 13 the binding one?
\end{exercise}

\begin{exercise}[$\star\star$]\label{exo:m1:margin:5}
From \cref{fig:m1:margin:offsets} give the saving from margining the
portfolio together, in dollars and in percent. A broker margins each leg
separately ``for prudence''. What does that prudence cost a client who pays
5\% a year for funding?
\end{exercise}

\begin{exercise}[$\star\star$]\label{exo:m1:margin:6}
A one-day 99\% value-at-risk is 2.1\% of a position. Give the margin for
margin periods of risk of two and five days under square-root scaling. Name
one reason why the true five-day figure could be higher and one why it could
be lower.
\end{exercise}

\begin{exercise}[$\star\star\star$]\label{exo:m1:margin:7}
\emph{Coding.} With \texttt{peak\_to\_trough} and \texttt{procyclical.csv}
reproduce the three ten-day increases quoted in
\cref{fig:m1:margin:procyclical} and the two calm-period averages. Find the
weight of the stressed floor for which the ten-day increase is 50\%, and its
cost in calm times.
\end{exercise}

\begin{exercise}[$\star\star\star$]\label{exo:m1:margin:8}
\emph{Find the flaw.} ``Our book is market neutral: long \$200 million of
shares at the \omterm{def:m1:financing:pb}{prime broker}, short \$200 million of index futures at the
clearing house. A 10\% fall in the market costs us nothing, so \omterm{def:m1:financing:margincall}{margin calls}
are not a risk for us.'' And a 10\% rise?
\end{exercise}

\section{Problem: A Fixed Portfolio in a Month Like March 2020}

\begin{problem}[{Weekend problem --- same contracts, more cash}]\label{pb:m1:margin:1}
The numbers are invented, in the proportions of a violent month. A fund is
long 400 index futures (multiplier \$50). At the start of the month the
future is at 3\,300 and the clearing house's \omterm{def:m1:clearing-and-settlement:margin}{initial margin} is 4.5\% of
notional. Four weeks later the future is at 2\,500 and the margin is 9\% of
notional. The fund keeps \$12 million of cash and holds the rest of its
assets in corporate bonds.

\medskip\noindent\textbf{Part I --- Two calls.}
\begin{enumerate}
\item Give the notional and the \omterm{def:m1:clearing-and-settlement:margin}{initial margin} at the start.
\item Give the \omterm{def:m1:clearing-and-settlement:margin}{variation margin} paid over the month.
\item Give the \omterm{def:m1:clearing-and-settlement:margin}{initial margin} at the end, and the change.
\item Give the total cash paid, in dollars and as a percentage of the
initial notional.
\item Compare with the fund's cash. What must it do?
\end{enumerate}

\medskip\noindent\textbf{Part II --- The worst day.}
On one day the future falls 9.5\% from 2\,750.
\begin{enumerate}[resume]
\item Give that day's \omterm{def:m1:clearing-and-settlement:margin}{variation margin}.
\item The clearing house makes an intraday call at noon, when the future is
down 6\%. Give the amount and the deadline's consequence.
\item The margin rate is raised that evening from 6\% to 7.5\% of notional.
Give the additional \omterm{def:m1:clearing-and-settlement:margin}{initial margin} at the new price.
\item Give the total cash called in twenty-four hours.
\end{enumerate}

\medskip\noindent\textbf{Part III --- Models.}
\begin{enumerate}[resume]
\item Using \cref{fig:m1:margin:procyclical}, which model would have called
the most additional margin in the first ten days, and which the least?
\item Which model leaves the clearing house best protected on day 15?
\item Explain why the floored model is unpopular in calm years.
\item A fund can choose between two brokers: one passes on the clearing
house's margin, the other charges 1.5 times it at all times. Which is safer
for the fund in a crisis?
\item The clearing house doubles margins. Show with one sentence of
arithmetic why the system as a whole needs more collateral although nobody
has become poorer.
\end{enumerate}

\medskip\noindent\textbf{Part IV --- Judgement.}
\begin{enumerate}[resume]
\item The fund's futures were a hedge of a short position in shares held
elsewhere. Does that help with the calls?
\item What should the fund have measured before the month began?
\item Why can the fund not simply post its corporate bonds?
\item If many funds sell bonds to meet calls, what happens next?
\item State the \emph{named result}: the additional \omterm{def:m1:clearing-and-settlement:margin}{initial margin} called on
the unchanged portfolio over the month.
\item In one sentence: what is the difference between being solvent and
being able to pay a \omterm{def:m1:financing:margincall}{margin call}?
\end{enumerate}
\end{problem}

\section{Interview questions}

\begin{interviewq}[$\star$ \iqroles{trader, researcher, developer}]\label{iq:m1:margin:1}
What is the difference between initial and \omterm{def:m1:clearing-and-settlement:margin}{variation margin}?
\end{interviewq}

\begin{interviewq}[$\star$ \iqroles{trader, bank}]\label{iq:m1:margin:2}
Describe how a scenario-based margin system such as SPAN computes the
requirement for a portfolio of futures and options.
\end{interviewq}

\begin{interviewq}[$\star\star$ \iqroles{researcher, bank}]\label{iq:m1:margin:3}
What is \omterm{def:m1:margin:procyclical}{procyclicality} of margin, and what tools reduce it? What do they
cost?
\end{interviewq}

\begin{interviewq}[$\star\star$ \iqroles{trader, researcher}]\label{iq:m1:margin:4}
You are short far out-of-the-money options. Which part of the margin
calculation is aimed at you, and why does it exist?
\end{interviewq}

\begin{interviewq}[$\star\star$ \iqroles{researcher, mle}]\label{iq:m1:margin:5}
How would you include margin in a backtest of a leveraged futures strategy?
\end{interviewq}

\begin{interviewq}[$\star\star\star$ \iqroles{trader, bank, researcher}]\label{iq:m1:margin:6}
Your firm is hedged across two clearing houses. Describe its liquidity risk
and how you would size the cash buffer.
\end{interviewq}
